% ============================================================
\chapter{Applications}
\label{ch:applications}
% ============================================================

Les th\'eor\`emes de points fixes ne sont pas des curiosit\'es abstraites~: ils sont les outils qui, dans les coulisses, font fonctionner une grande partie de l'analyse classique. Le th\'eor\`eme d'existence et d'unicit\'e de Picard-Lindel\"of pour les \'equations diff\'erentielles ordinaires~? C'est le principe de contraction de Banach appliqu\'e \`a un op\'erateur int\'egral. Les \'equations int\'egrales de Fredholm~? Encore Banach. Le th\'eor\`eme des fonctions implicites, pilier du calcul diff\'erentiel~? Le principe de contraction, une fois de plus. Ce chapitre rassemble ces applications classiques, montrant comment un seul r\'esultat abstrait --- l'existence d'un point fixe pour une contraction --- irrigue des pans entiers des math\'ematiques appliqu\'ees.

\begin{intuition}
Ce chapitre illustre la puissance des th\'eor\`emes de points fixes en les appliquant \`a
trois domaines classiques de l'analyse~: les \'equations diff\'erentielles ordinaires (via
le th\'eor\`eme de Picard-Lindel\"of comme cons\'equence du principe de contraction de
Banach), les \'equations int\'egrales (Fredholm et Volterra), et le th\'eor\`eme des
fonctions implicites. Dans chaque cas, le probl\`eme est reformul\'e comme un probl\`eme de
point fixe dans un espace fonctionnel ad\'equat.
\end{intuition}

% ------------------------------------------------------------
\section{\'Equations diff\'erentielles ordinaires~: Picard-Lindel\"of}
% ------------------------------------------------------------

\begin{theoreme}[Picard-Lindel\"of (Cauchy-Lipschitz)]
\label{thm:picard-lindelof}
\index{th\'eor\`eme de Picard-Lindel\"of}
Soit $f\colon [t_0 - a, t_0 + a] \times \overline{B}(y_0, b) \to \R^n$ une application
continue, lipschitzienne en $y$ de constante $L$~:
\[
  \norm{f(t, y_1) - f(t, y_2)} \leq L \norm{y_1 - y_2}
\]
pour tous $t, y_1, y_2$. Posons $M = \sup \norm{f}$ et $\delta = \min(a, b/M)$.
Alors le probl\`eme de Cauchy
\[
  y'(t) = f(t, y(t)), \qquad y(t_0) = y_0
\]
admet une unique solution $y \in C^1([t_0 - \delta, t_0 + \delta], \R^n)$.
\end{theoreme}

\begin{proof}
Le probl\`eme de Cauchy est \'equivalent \`a l'\'equation int\'egrale
\[
  y(t) = y_0 + \int_{t_0}^{t} f(s, y(s)) \, ds =: (Ty)(t).
\]
Consid\'erons l'espace $E = \{y \in C([t_0 - \delta, t_0 + \delta], \R^n) :
\norm{y(t) - y_0} \leq b \; \forall t\}$ muni de la norme
\[
  \norm{y}_\lambda = \sup_{t} e^{-\lambda \abs{t - t_0}} \norm{y(t) - y_0}
\]
pour $\lambda > L$ \`a choisir.

\emph{$T$ envoie $E$ dans $E$}~: Pour $y \in E$,
$\norm{(Ty)(t) - y_0} \leq M \abs{t - t_0} \leq M\delta \leq b$.

\emph{$T$ est une contraction}~: Pour $y_1, y_2 \in E$,
\begin{align*}
  e^{-\lambda\abs{t-t_0}} \norm{(Ty_1)(t) - (Ty_2)(t)}
  &\leq e^{-\lambda\abs{t-t_0}} \int_{t_0}^{t} L \norm{y_1(s) - y_2(s)} \, ds \\
  &\leq L \norm{y_1 - y_2}_\lambda \int_{t_0}^{t} e^{\lambda(\abs{s-t_0} - \abs{t-t_0})} ds \\
  &\leq \frac{L}{\lambda} \norm{y_1 - y_2}_\lambda.
\end{align*}
En choisissant $\lambda > L$, on obtient $\norm{Ty_1 - Ty_2}_\lambda \leq
\frac{L}{\lambda} \norm{y_1 - y_2}_\lambda$ avec $L/\lambda < 1$.

Par le principe de contraction de Banach, $T$ admet un unique point fixe dans $E$.
\end{proof}

\begin{remarque}
La norme $\norm{\cdot}_\lambda$ (norme de Bielecki) est une astuce classique qui \'evite de
devoir restreindre l'intervalle de temps pour obtenir la contraction.
\end{remarque}

\begin{corollaire}[D\'ependance continue]
La solution $y(t; t_0, y_0)$ d\'epend contin\^ument des donn\'ees initiales $(t_0, y_0)$.
\end{corollaire}

% ------------------------------------------------------------
\section{\'Equations int\'egrales de Fredholm}
% ------------------------------------------------------------

\begin{definition}[\'Equation de Fredholm de seconde esp\`ece]
\index{\'equation de Fredholm}
Soient $K \in C([a,b]^2, \R)$ un \emph{noyau}, $f \in C([a,b], \R)$ et
$\lambda \in \R$. L'\emph{\'equation de Fredholm de seconde esp\`ece} est~:
\[
  u(t) = f(t) + \lambda \int_a^b K(t,s) \, u(s) \, ds.
\]
\end{definition}

\begin{theoreme}[Existence et unicit\'e pour Fredholm]
\label{thm:fredholm}
Si $\abs{\lambda} < \frac{1}{(b-a) \norm{K}_\infty}$, alors l'\'equation de Fredholm
admet une unique solution $u \in C([a,b], \R)$.
\end{theoreme}

\begin{proof}
D\'efinissons l'op\'erateur $T\colon C([a,b]) \to C([a,b])$ par
\[
  (Tu)(t) = f(t) + \lambda \int_a^b K(t,s) \, u(s) \, ds.
\]
Alors~:
\[
  \norm{Tu_1 - Tu_2}_\infty \leq \abs{\lambda} (b-a) \norm{K}_\infty \norm{u_1 - u_2}_\infty.
\]
Sous l'hypoth\`ese $\abs{\lambda} (b-a) \norm{K}_\infty < 1$, l'op\'erateur $T$ est une
contraction sur l'espace de Banach $C([a,b])$, et le th\'eor\`eme de Banach s'applique.
\end{proof}

\begin{exemple}
L'\'equation $u(t) = \sin(t) + \frac{1}{4} \int_0^1 ts \, u(s) \, ds$ a un noyau
$K(t,s) = ts$ avec $\norm{K}_\infty = 1$. Comme $\abs{\lambda}(b-a)\norm{K}_\infty
= \frac{1}{4} < 1$, la solution existe et est unique.
\end{exemple}

% ------------------------------------------------------------
\section{\'Equations int\'egrales de Volterra}
% ------------------------------------------------------------

\begin{definition}[\'Equation de Volterra de seconde esp\`ece]
\index{\'equation de Volterra}
L'\emph{\'equation de Volterra de seconde esp\`ece} est~:
\[
  u(t) = f(t) + \lambda \int_a^t K(t,s) \, u(s) \, ds.
\]
La diff\'erence avec Fredholm est que la borne sup\'erieure d'int\'egration est $t$
(et non $b$).
\end{definition}

\begin{theoreme}[Existence et unicit\'e pour Volterra]
\label{thm:volterra}
Pour tout $\lambda \in \R$ et tout noyau $K$ continu, l'\'equation de Volterra de seconde
esp\`ece admet une unique solution $u \in C([a,b], \R)$.
\end{theoreme}

\begin{proof}
L'op\'erateur $T$ d\'efini par $(Tu)(t) = f(t) + \lambda \int_a^t K(t,s) u(s) ds$ n'est
pas n\'ecessairement une contraction pour la norme uniforme. Cependant, utilisons la norme
de Bielecki~:
\[
  \norm{u}_\mu = \sup_{t \in [a,b]} e^{-\mu(t-a)} \abs{u(t)}.
\]
Alors~:
\begin{align*}
  e^{-\mu(t-a)} \abs{(Tu_1)(t) - (Tu_2)(t)}
  &\leq \abs{\lambda} \norm{K}_\infty \int_a^t e^{-\mu(t-s)} e^{-\mu(s-a)} \abs{u_1(s) - u_2(s)} \, ds \\
  &\leq \abs{\lambda} \norm{K}_\infty \norm{u_1 - u_2}_\mu \int_a^t e^{-\mu(t-s)} ds \\
  &\leq \frac{\abs{\lambda} \norm{K}_\infty}{\mu} \norm{u_1 - u_2}_\mu.
\end{align*}
En choisissant $\mu > \abs{\lambda} \norm{K}_\infty$, $T$ est une contraction.
\end{proof}

\begin{attention}
Contrairement \`a Fredholm, l'\'equation de Volterra admet toujours une solution unique,
sans condition sur $\lambda$. C'est une cons\'equence du caract\`ere \og{}triangulaire\fg
de l'op\'erateur de Volterra.
\end{attention}

% ------------------------------------------------------------
\section{Th\'eor\`eme des fonctions implicites}
% ------------------------------------------------------------

\begin{theoreme}[Fonctions implicites via point fixe]
\label{thm:implicit}
\index{th\'eor\`eme des fonctions implicites}
Soit $F\colon U \to \R^m$ de classe $C^1$ o\`u $U$ est un ouvert de $\R^n \times \R^m$.
Soit $(a, b) \in U$ tel que $F(a, b) = 0$ et $D_y F(a, b)$ est inversible. Alors il
existe des voisinages $V$ de $a$ et $W$ de $b$ et une unique application
$\varphi\colon V \to W$ de classe $C^1$ telle que
\[
  F(x, \varphi(x)) = 0 \quad \forall x \in V.
\]
\end{theoreme}

\begin{proof}
On ram\`ene le probl\`eme \`a un point fixe. Posons $A = D_y F(a,b)$ et
$G(x,y) = y - A^{-1} F(x,y)$. Alors $F(x,y) = 0$ \'equivaut \`a $y = G(x,y)$.

Montrons que pour $x$ proche de $a$, $y \mapsto G(x,y)$ est une contraction au voisinage
de $b$. On a $D_y G(a,b) = I - A^{-1} D_y F(a,b) = 0$. Par continuit\'e de $D_y G$, il
existe $r > 0$ tel que $\norm{D_y G(x,y)} \leq 1/2$ pour $(x,y)$ dans un voisinage de
$(a,b)$.

Ainsi, pour $x$ fix\'e proche de $a$, $y \mapsto G(x,y)$ est une contraction de constante
$1/2$ sur $\overline{B}(b, r)$ (quitte \`a ajuster $r$). Par le th\'eor\`eme de Banach,
il existe un unique $y = \varphi(x)$ tel que $G(x, \varphi(x)) = \varphi(x)$, soit
$F(x, \varphi(x)) = 0$.

La r\'egularit\'e $C^1$ de $\varphi$ s'obtient par diff\'erentiation de la relation
$F(x, \varphi(x)) = 0$, ce qui donne
$D\varphi(x) = -[D_y F(x, \varphi(x))]^{-1} D_x F(x, \varphi(x))$.
\end{proof}

\begin{remarque}
Cette d\'emonstration par point fixe est plus constructive que la d\'emonstration classique
(par le th\'eor\`eme d'inversion locale) car elle fournit un algorithme d'approximation~:
la suite $y_{n+1} = G(x, y_n)$ converge vers $\varphi(x)$.
\end{remarque}

% ------------------------------------------------------------
\section{Compl\'ements~: m\'ethode de Newton et points fixes}
% ------------------------------------------------------------

\begin{proposition}[Convergence de Newton comme point fixe]
La m\'ethode de Newton pour r\'esoudre $F(x) = 0$ dans $\R^n$ s'\'ecrit~:
\[
  x_{n+1} = x_n - [DF(x_n)]^{-1} F(x_n) =: T(x_n).
\]
Sous des conditions de r\'egularit\'e (Lipschitz sur $DF$, inversibilit\'e de $DF(x^*)$),
l'application $T$ n'est pas une contraction globale, mais elle l'est localement au
voisinage de la solution $x^*$, avec un taux de convergence quadratique~:
$\norm{x_{n+1} - x^*} \leq C \norm{x_n - x^*}^2$.
\end{proposition}

\begin{formules}[title=\textbf{R\'ecapitulatif des applications}]
\begin{center}
\renewcommand{\arraystretch}{1.3}
\begin{tabular}{|l|l|l|}
\hline
\textbf{Probl\`eme} & \textbf{Th\'eor\`eme utilis\'e} & \textbf{Espace fonctionnel} \\
\hline
EDO (Picard-Lindel\"of) & Banach & $C([t_0-\delta, t_0+\delta], \R^n)$ \\
Fredholm & Banach & $C([a,b], \R)$ \\
Volterra & Banach (norme Bielecki) & $C([a,b], \R)$ \\
Fonctions implicites & Banach & $\R^m$ \\
\hline
\end{tabular}
\end{center}
\end{formules}

% ------------------------------------------------------------
\section{Exercices}
% ------------------------------------------------------------

\begin{exercice}
Consid\'erons l'EDO $y' = y^2$, $y(0) = 1$. Calculer les trois premi\`eres it\'er\'ees de
Picard $y_0(t) = 1$, $y_{n+1}(t) = 1 + \int_0^t y_n(s)^2 ds$. Comparer avec la solution
exacte $y(t) = \frac{1}{1-t}$.
\end{exercice}

\begin{exercice}
R\'esoudre l'\'equation de Fredholm $u(t) = 1 + \lambda \int_0^1 u(s) ds$ explicitement
pour tout $\lambda \neq 1$. V\'erifier que la solution diverge quand $\lambda \to 1$.
\end{exercice}

\begin{exercice}
R\'esoudre l'\'equation de Volterra $u(t) = 1 + \int_0^t u(s) ds$ en reconnaissant une
EDO. V\'erifier le r\'esultat par la m\'ethode des it\'er\'ees successives.
\end{exercice}

\begin{exercice}
Montrer que l'\'equation $F(x, y) = x^2 + y^2 - 1 = 0$ d\'efinit $y$ comme fonction de
$x$ au voisinage de $(0, 1)$. Calculer $\varphi'(0)$ et $\varphi''(0)$.
\end{exercice}

\begin{exercice}
Soit $K(t,s) = e^{ts}$ et $f(t) = 1$. \'Etudier la convergence des it\'er\'ees de
l'op\'erateur $T$ associ\'e \`a l'\'equation de Fredholm
$u(t) = 1 + \lambda \int_0^1 e^{ts} u(s) ds$ en fonction de $\lambda$.
\end{exercice}

\begin{exercice}[Syst\`emes d'EDO]
Appliquer le th\'eor\`eme de Picard-Lindel\"of au syst\`eme $y_1' = y_2$, $y_2' = -y_1$,
$y_1(0) = 1$, $y_2(0) = 0$. Calculer explicitement les it\'er\'ees de Picard et montrer
qu'elles convergent vers $(\cos t, -\sin t)$.
\end{exercice}
